%% Independent methods file. Not woven into main.tex prose.
%% Every quantity described here resolves to a hash-bound evidence entry.

\subsection{The corpus and the frozen split}

The locally held corpus is the Indiana University chest X-ray collection
distributed through Open-i~\cite{demnerfushman2016preparing}, which pairs
de-identified radiographs with their free-text reports. Reports, not images, are
the unit of everything reported here.

The split was frozen before any of the analyses in this paper. It is
patient-disjoint by construction: a report's assignment is a function of a
salted hash of its patient identifier, so no patient contributes to more than one
side and the assignment can be recomputed by anyone holding the same identifier
list. The frozen inventory holds \NPatients{} patients and \NImages{} images,
and divides into \NTrain{} training, \NVal{} validation and \NTest{} test
reports.

One property of that inventory is worth stating because it is the kind of thing
usually smoothed over. A total of \NOutsideInventory{} radiograph files on disk
are not in
the inventory the split was built from. The project's recorded decision was to
leave them visible and outside rather than manufacture inventory rows to absorb
them, which keeps the split reproducible from the inventory alone and keeps the
discrepancy countable. This paper reports the discrepancy rather than the larger
file count.

\subsection{The pre-registered subsample}

Labelling was applied to a subsample of \NCohort{} reports from the frozen test
split, which is \CoveragePct{} per cent of it. The subsample is not the first
\NCohort{} identifiers in any ordering that correlates with content: the
selection rule sorts test identifiers by a salted hash and takes the first
\NCohort{}, with the salt fixed and recorded before the labels were read. The
recorded contract states explicitly that the result is not a lexicographic
prefix.

This matters for interpreting one of the results below. A pre-registered
selection rule protects against choosing a sample after seeing its labels. It
does not protect against the sample being unusual, and the two are routinely
conflated.

\subsection{Labels}

Labels are the \NFindings{} observations of the CheXpert schema~\cite{irvin2019chexpert},
assigned to report text by the released CheXbert labeller~\cite{smit2020chexbert},
which is the labeller the standard automatic metric in this area is built on.
The checkpoint used is recorded with its SHA-256 checksum in the recorded artifact, so
the labelling can be tied to a specific set of weights rather than to a name.

The labeller emits four values per observation: positive, negative, uncertain,
and not mentioned. Throughout this paper an observation counts as a gold positive
only when the labeller assigns the positive value. Uncertain and unmentioned are
not counted as positives, and they are not counted as negatives either when a
count of positives is what is being reported.

The unmentioned value dominates and its interpretation is a genuine caveat.
Across the \NFindings{} observations the median count of test reports that do not
mention an observation at all is \MedianBlankPct{} per cent of the split, and for
\MaxBlankFinding{} it is \MaxBlank{} of \NTest{} reports. A report that is silent
about an observation is not a report asserting its absence, and any analysis that
treats the two as equivalent is making a convention visible only in its methods
section. We avoid the convention by reporting positives and misses on positives
rather than accuracy over all four values.

For the four-state composition figure, let \(n_{f,s}\) be the count of state
\(s\in\{+, -, ?, \mathrm{blank}\}\) for finding \(f\) in the frozen test split.
The build checks \(\sum_s n_{f,s}=\NTest{}\) for every finding and plots the
fractions \(n_{f,s}/\NTest{}\). The blank state is therefore retained as a
separate denominator component and is never interpreted as a negative label;
the resulting display describes CheXbert output rather than clinical prevalence.

\subsection{The presence gate and the two power questions}

The analysis applies a presence gate fixed before the runs: an observation is
scored only if the subsample holds at least \MinPresent{} gold positives for it,
and otherwise its cell is reported as underpowered with no rate attached. The
gate value is not asserted here but recovered: the recorded per-observation
status of both control arms is checked against the positive counts, and the build
stops if any status disagrees with a gate of \MinPresent{}.

We keep the eligibility sets explicit. Let \(S\) be the frozen test reports,
\(C\subseteq S\) the \NCohort{}-report draw, and \(P_f=\{i\in S:
y_{if}=+\}\) the reports labelled positive for finding \(f\). The pilot
denominator is
\[
  K_f = |C\cap P_f|,
  \qquad
  \mathrm{status}(f)=
  \begin{cases}
    \mathrm{scored}, & K_f\geq \MinPresent{},\\
    \mathrm{underpowered}, & K_f<\MinPresent{}.
  \end{cases}
\]
At the benchmark level, a cell is eligible only when the permission,
artifact-availability, and no-substitution gates all hold; those gates are
recorded separately from \(K_f\). Thus an inaccessible cell, an observable but
underpowered cell, and a scored cell have different denominators and are never
collapsed into one missing-value code.

Two distinct power questions follow, and separating them is one of this paper's
points.

The first asks what the corpus can support. For an observation present in
\(K\) of the \(N\) test reports, a draw of \(n\) reports contains \(nK/N\)
positives in expectation, so the smallest draw whose expectation reaches the gate
is \(\lceil \text{\MinPresent{}} \cdot N / K \rceil\). When that number exceeds
\(N\) itself, no subsample of the test split can reach the gate and the
limitation is the corpus, not the pilot. This is arithmetic on counts, not a
simulation.

The second asks whether the draw that was actually taken resembles the split it
came from. For each observation we compute the exact two-sided hypergeometric
probability of the observed positive count, summing the probability of every
outcome no more likely than the one observed, under the null that the draw is
uniform over the split. With \NFindings{} observations tested the Bonferroni
threshold is \Bonferroni{}, and we report probabilities against that threshold
rather than against \(0.05\).

\subsection{The two controls}

Two controls are scored in label space on the subsample. The majority control
emits, for every report, the single most frequent findings text in the training
split; it is constant by construction, which the analysis verifies by requiring
each of its per-observation positive counts to be either zero or all \NCohort{}
rows. The retrieval control embeds the indication field with term-frequency--inverse-document-frequency (TF--IDF) weighting, retrieves
the nearest training report by cosine similarity, and returns that report's
findings text; retrieval is restricted to the training side, so a report can
never retrieve itself or a study from its own patient.

For each observation and each control we count the gold positives the control
failed to assert, and report that count as a rate with a one-sided \(90\%\)
Clopper--Pearson upper limit~\cite{clopper1934use}. The upper limit is reported
because the counts are small: at these sample sizes an interval is the result and
a point estimate on its own would misrepresent it.

For a scored finding \(f\), the reported miss rate is
\(\widehat q_f=M_f/K_f\), where \(M_f\) is the number of positive reports the
control failed to assert. The denominator is therefore the observed
gold-positive count in \(C\), not \(|C|\), \(|S|\), or the number of files on
disk. No \(\widehat q_f\) is defined when \(K_f<\MinPresent{}\).

The source design asked for a coverage--risk curve over abstention rates.
The bound control tables record a single coverage of the labelled draw and
no ranking that would order reports for intermediate abstention. Only the
\NCoverageRiskPoints{} cells with \(K_f\geq\MinPresent{}\) can contribute a
point, and each point is \(\widehat q_f\) at coverage
\CoverageRiskCoverage{}. A curve is therefore not drawn. The same tables
record the contact-gated evaluation as \TypeRBound{}; no score is computed
or printed under that name.

The same two controls also exist scored a different way, on surface text over the
whole test split rather than in label space on the subsample, using exact match
and unigram token F1. Those numbers are reported here only as a contrast, and
they are not comparable to the label-space numbers.

\subsection{Quantities excluded by design}

Three quantities were within the project's reach and were kept out of the
analysis by design, because each would present the result as stronger than
the access situation supports.

Scores are reported under the name of the corpus they were computed on. The
locally held corpus is a different corpus from the one the benchmark scores,
and a number computed on it carries that corpus's name rather than the
benchmark's, since the substitution would be a category error rather than an
approximation. The project's own scoring wrapper enforces this by accepting
only samples of the benchmark's split size.

The label counts are reported as counts of positives and misses. They are kept
out of the field's headline automatic report-generation aggregate, because
that conversion would invite comparison against published values computed on a
different corpus with a different sample size.

The one local inference run is reported as what it is. A run of a
chest-radiograph foundation model over \HarnessImages{} images exists and is
recorded as a smoke test that checks the code runs end to end, explicitly
marked as not the official evaluation. It contributes no number to this paper.

\subsection{Evidence discipline}

No number in this manuscript is typed by hand. Every artifact it rests on is
recorded with its SHA-256 checksum in a manifest that accompanies the manuscript
source; the figure and table code verifies each checksum against the bytes on disk
before reading them and aborts on any mismatch, and it emits both the figures and
the numeric macros the text prints. Several consistency checks are enforced at
build time rather than asserted in prose: the two control arms must agree on how
many gold positives the subsample holds, the recomputed list of observations
expected to yield fewer than one positive must match the list the runs recorded,
the presence gate must be recoverable from the recorded statuses, and the set of
access obstacles the first figure describes must equal the set the project
recorded. Any disagreement stops the build instead of producing a manuscript.
